\documentclass[11pt,a4paper]{article}
\input{preambolo}

\title{\bfseries Simulazione quantistica di un trimero di spin ad anello:\\
stato fondamentale, dinamica e correlatori,\\
in assenza e in presenza di rumore}
\author{Samuele Schiavi \\ \small Tesi triennale in Fisica, Università di Parma
--- Relatore: Prof.\ Alessandro Chiesa}
\date{}

\begin{document}
\maketitle
\thispagestyle{fancy}

\begin{abstract}
\noindent
Si studia un sistema di tre spin $1/2$ accoppiati ciclicamente (trimero ad
anello), estensione del caso più semplice del dimero già trattato in
precedenza, con l'obiettivo di simulare quantisticamente le sue proprietà
statiche e dinamiche tramite il Variational Quantum Eigensolver (VQE) e la
decomposizione di Suzuki--Trotter. Il lavoro è organizzato in due parti.
Nella prima si tratta il sistema come chiuso: si determina lo spettro
esatto, si costruisce un ansatz variazionale per lo stato fondamentale, si
implementa l'evoluzione temporale reale e si calcolano correlatori
dinamici a due tempi, tutto su simulatore a vettore di stato ideale. Nella
seconda si introduce un modello di rumore realistico, calibrato su
hardware IBM, e si ripetono le stesse misure nel formalismo dell'operatore
densità, per quantificare quanto ciascuno strumento sopravvive
all'imperfezione dell'hardware. Rispetto al dimero, il trimero introduce
un ingrediente fisico assente nel caso a due spin --- la frustrazione
geometrica --- insieme a complicazioni tecniche concrete che ne derivano:
più forme possibili per il termine di interazione antisimmetrica, un
ansatz variazionale a più parametri il cui confronto capovolge la
conclusione già stabilita sul dimero, e una struttura di Trotter a più
livelli annidati.
\end{abstract}

\tableofcontents

% =====================================================================
\part*{Parte I --- Il sistema chiuso}
\addcontentsline{toc}{part}{Parte I --- Il sistema chiuso}
% =====================================================================

\section{Il modello fisico}
\label{sec:modello}

\subsection{Hamiltoniana e topologia}

Il sistema è costituito da tre spin $1/2$ disposti ai vertici di un
triangolo isoscele: un legame di scambio $J$ chiude la base, fra i siti
$1$ e $2$, e due legami $J'$ ne formano i lati, fra le coppie $(2,3)$ e
$(3,1)$. Il caso equilatero frustrato corrisponde a $J'=J$; nel presente
lavoro si mantiene fisso il rapporto $J'/J=0.4$, con $J=1$ come unità di
energia. In presenza di un campo assiale $b$ lungo $z$ e di un termine di
interazione antisimmetrica di Dzyaloshinskii--Moriya (DM) di intensità
$D$, l'Hamiltoniana completa è
\begin{equation}
H = \underbrace{J\,\vec\sigma_1\!\cdot\!\vec\sigma_2 +
J'\,(\vec\sigma_2\!\cdot\!\vec\sigma_3+\vec\sigma_3\!\cdot\!\vec\sigma_1)}_{\text{scambio isotropo}}
+ \underbrace{b\,(Z_1+Z_2+Z_3)}_{\text{campo}}
+ \underbrace{D\sum_{\langle i,j\rangle} (X_iZ_j-Z_iX_j)}_{\text{termine DM}},
\label{eq:H}
\end{equation}
con $\vec\sigma_i\cdot\vec\sigma_j=X_iX_j+Y_iY_j+Z_iZ_j$. Come nel caso del
dimero, ogni spin $1/2$ si mappa direttamente su un qubit fisico, senza
alcun mapping fermionico intermedio: il costo di rappresentazione non
cambia passando da due a tre siti. Il legame $J'$ fra $(3,1)$, assente in
una catena aperta a tre siti, è ciò che chiude il ciclo e genera la
frustrazione discussa di seguito.

\subsection{Frustrazione geometrica}

Con scambio antiferromagnetico ($J,J'>0$), ogni legame favorisce
energeticamente una configurazione con i due spin antiparalleli. Su un
ciclo di lunghezza pari questa condizione può essere soddisfatta su tutti
i legami simultaneamente (configurazione alternata). Su un ciclo di
lunghezza \emph{dispari}, come il triangolo qui considerato, questo non è
mai possibile: comunque si orientino i tre spin, almeno un legame resta
con i suoi due spin paralleli, energeticamente sfavorito. Questa è la
frustrazione geometrica: un fatto strutturale, che dipende dalla parità
del ciclo e non dai valori specifici di $J,J'$. Un sistema a due spin (il
dimero) non può mostrare questo fenomeno per costruzione, avendo un solo
legame.

\subsection{Spettro esatto: decomposizione di Kambe}

La simmetria di scambio fra i due siti della base, $1\leftrightarrow2$,
permette di risolvere lo spettro in forma chiusa componendo prima i
momenti angolari della coppia $(1,2)$ in uno spin totale $S_{12}\in\{0,1\}$
e poi accoppiando il terzo spin al risultato (decomposizione di Kambe).
Per $D=0$ si ottengono tre multipletti:
\begin{align}
\text{A}\ (S_{12}{=}1,\,S{=}3/2):&\quad E_A = J+2J', \\
\text{B}\ (S_{12}{=}1,\,S{=}1/2):&\quad E_B = J-4J', \\
\text{C}\ (S_{12}{=}0,\,S{=}1/2):&\quad E_C = -3J,
\end{align}
ciascuno separato in campo secondo $E(S,M)=E_\text{blocco}+2bM$. Il campo
critico, al quale il fondamentale cambia identità, dipende da quale
multipletto sia il fondamentale a $b=0$:
\begin{equation}
b_c = \begin{cases} 2J+J' & \text{se il fondamentale a } b{=}0 \text{ è C}, \\
3J' & \text{se è B.}\end{cases}
\end{equation}
Al punto di lavoro adottato in questo capitolo ($J'/J=0.4$), il
fondamentale a $b=0$ è il blocco C, e quindi $b_c=2J+J'=2.4\,J$. Lo
spettro numerico, ottenuto per diagonalizzazione diretta della matrice
$8\times8$, coincide con queste formule analitiche a precisione macchina
su tutta la griglia di campo esplorata.

\begin{figure}[htbp]
\centering
\includegraphics[width=0.78\textwidth]{figure/fig01_spettro_anello}
\caption{Spettro del trimero isoscele ($J'/J=0.4$), $D=0$, in funzione del
campo $b/J$. Otto livelli in tre famiglie secondo la decomposizione di
Kambe (blocchi A, B, C). La curva che segue lo stato fondamentale parte
dal blocco C e incrocia il blocco A al campo critico $b_c/J=2.4$.}
\label{fig:spettro}
\end{figure}
\FloatBarrier

Senza termine DM, l'Hamiltoniana conserva sia la magnetizzazione totale
$M$ sia il quadrato del momento angolare della coppia base, $S_{12}^2$:
due numeri quantici buoni, non uno solo. L'incrocio a $b_c$ avviene fra
il blocco C ($S_{12}=0$) e il blocco A ($S_{12}=1$): stati con $S_{12}^2$
diverso non si mescolano finché questa simmetria resta intatta, e
l'incrocio resta quindi un incrocio vero, non evitato.

\subsection{Il termine di Dzyaloshinskii--Moriya: scelta della forma}
\label{sec:dm}

A differenza del dimero, dove il termine DM ha un'unica forma possibile
(un solo legame), con tre legami esistono più modi di accendere
un'interazione antisimmetrica, e non tutti rompono l'incrocio protetto
discusso sopra. L'intera costruzione teorica del trimero isoscele
presuppone l'invarianza per scambio dei siti $1,2$, rappresentata
dall'operatore di permutazione $P_{12}$: un termine DM aggiunto senza
vincoli rischierebbe di rompere questa simmetria in un modo non
giustificato dalla geometria del problema.

Sono state verificate sistematicamente tutte le $12$ combinazioni di
segno di $(D_{12},D_{23},D_{31})\propto(J,J',J')$ (fattore di intensità
comune $r$), calcolando $\max\lVert P_{12}H_\text{DM}P_{12}^\dagger -
H_\text{DM}\rVert$ per ciascuna. Il risultato è netto: una sola
disposizione di segno (a meno dello scambio $23\leftrightarrow31$)
rispetta esattamente la simmetria.

\begin{table}[htbp]
\centering
\renewcommand{\arraystretch}{1.4}
\begin{tabular}{@{}p{2.3cm}p{6.3cm}p{6.3cm}@{}}
\toprule
\rowcolor{grigetto}
& \textbf{Opzione A} (simmetrica) & \textbf{Opzione B} (adottata) \\
\midrule
forma & $D'\big[(X_2Z_3{-}Z_2X_3)-(X_3Z_1{-}Z_3X_1)\big]$, $D_{12}=0$ &
$D(X_1Z_2{-}Z_1X_2)+0.4D(X_2Z_3{-}Z_2X_3)+0.4D(X_3Z_1{-}Z_3X_1)$ \\
simmetria $P_{12}$ & rispettata esattamente & rotta \\
$\lVert[H,S_{12}^2]\rVert$ & $\sim10^{-16}$ (esatto) & $\sim1.4$ a
$D/J{=}0.15$ \\
struttura residua & blocco $2{\times}2$ (C) $+$ blocco $6{\times}6$
(A$+$B) & nessuna, $8\times8$ pieno \\
apre l'incrocio? & no & sì \\
\bottomrule
\end{tabular}
\caption{Le due forme del termine DM sopravvissute al vincolo di
simmetria (o alla sua assenza). $D$ è il parametro di intensità
complessivo; i coefficienti $0.4$ nell'Opzione B derivano dal rapporto
$J'/J=0.4$ fissato in questo lavoro.}
\label{tab:dm-opzioni}
\end{table}
\FloatBarrier

La ragione per cui l'incrocio resta protetto sotto l'Opzione A, anche a
intensità $D$ arbitraria, è la regola di non-incrocio di
von~Neumann--Wigner: poiché l'Opzione A commuta con $S_{12}^2$,
l'Hamiltoniana resta a blocchi nei due settori $S_{12}=0$ e $S_{12}=1$, e
l'elemento di matrice fra uno stato del blocco C e uno stato del blocco A
è esattamente nullo per ogni valore di $b$. Due livelli senza
accoppiamento diretto possono incrociarsi liberamente al variare di un
solo parametro, senza necessità di alcun aggiustamento fine. Sotto
l'Opzione B, che non commuta con $S_{12}^2$, quell'elemento di matrice è
generalmente non nullo: la regola non si applica più, e l'incrocio
diventa evitato, con un gap minimo genuino.

Numericamente: il gap minimo cercato in un intorno di $b_c$ (non a $b_c$
fisso) vale $2.3\times10^{-8}$ a $b/J=2.372$ per l'Opzione A --- zero a
precisione numerica, l'incrocio resta chiuso qualunque sia l'intensità
del DM in questa forma --- e $0.2543$ a $b/J=2.407$ per l'Opzione B. Che
l'incrocio resti vero sotto l'Opzione A non significa che resti fermo:
sotto questa opzione l'Hamiltoniana si scompone esattamente in un blocco
$2\times2$ isolato (settore C, accoppiamento con il resto nullo a
precisione macchina) e un blocco $6\times6$ pieno (settore $S_{12}=1$,
dove i sotto-blocchi A e B si mescolano fra loro, con elementi fuori
diagonale fino a $1.13$). È questo mescolamento interno al settore
$S_{12}=1$, non un mescolamento con il settore C, a spostare il ramo più
basso di quel settore e con esso il punto di incrocio, da $b/J=2.4$ a
$b/J\approx2.372$; il blocco C, isolato, resta esattamente sulla formula
$D=0$ per costruzione.

\begin{figure}[htbp]
\centering
\includegraphics[width=\textwidth]{figure/fig02_gap_magnetizzazione_anello}
\caption{\textbf{Sinistra:} distanza fondamentale--primo eccitato, $D=0$
contro le due opzioni di DM ($D/J=0.15$). L'Opzione A resta
indistinguibile da $D=0$ esattamente al punto critico; l'Opzione B apre
un gap minimo vero, $0.254$, spostato a $b/J=2.407$.
\textbf{Destra:} magnetizzazione dello stato fondamentale, $D=0$ contro
Opzione B: il salto netto diventa un crossover continuo.}
\label{fig:gap}
\end{figure}
\FloatBarrier

Nel seguito di questo lavoro si adotta sempre l'\textbf{Opzione B}: è
l'unica scelta di DM per cui il problema variazionale al punto critico è
ben posto (l'incrocio protetto dall'Opzione A lascerebbe il fondamentale
degenere, senza un bersaglio unico per il VQE) e per cui il trimero
riproduce, punto per punto, lo stesso ruolo che il termine DM ha nel
dimero.

\section{Ricerca dello stato fondamentale: VQE}
\label{sec:vqe}

\subsection{Ansatz a confronto}

Il VQE cerca $\boldsymbol\theta^*=\arg\min_{\boldsymbol\theta}
\bra{\psi(\boldsymbol\theta)}H\ket{\psi(\boldsymbol\theta)}$. Poiché
l'energia esatta è nota dalla decomposizione di Kambe (\S\ref{sec:modello}),
il VQE viene qui usato come strumento diagnostico per misurare quanto
bene un dato circuito rappresenti lo stato fondamentale vero, tramite la
fidelity $\mathcal F=\lvert\langle\psi_\text{VQE}|\psi_0\rangle\rvert^2$.

Sono state confrontate due famiglie di ansatz, con la stessa struttura
esterna: una rotazione $X$ sul primo qubit, seguita da un blocco a due
qubit per ciascuno dei tre legami (tre parametri), seguita da tre
rotazioni $R_y$ indipendenti, una per qubit (altri tre parametri) --- sei
parametri in totale per entrambe le famiglie. Le due famiglie differiscono
solo nel blocco a due qubit usato sui legami:
\begin{itemize}[leftmargin=1.6em,itemsep=2pt]
\item \textbf{RBS}$(\varphi)$: una rotazione di Givens reale, realizzata
con due CZ attorno a una coppia di rotazioni $R_y(\varphi),R_y(-\varphi)$,
racchiuse fra porte di Hadamard --- un solo parametro libero $\varphi$
nonostante i due gate $R_y$;
\item \textbf{$W(\theta)$}: due CNOT attorno a un'unica $R_y(\theta)$, la
famiglia di ansatz canonico di riferimento (Crippa et al., 2021).
\end{itemize}

Senza termine DM, al punto critico $b_c/J=2.4$, entrambe le famiglie
raggiungono $\mathcal F=1$ a precisione numerica: nessuna differenza
osservabile, un caso non discriminante fra le due scelte.

\subsection{Confronto sotto il termine DM}
\label{sec:rbs-vs-w}

Il confronto diventa discriminante al punto di lavoro denominato
\emph{Scenario A} ($b/J=b_c=2.4$, $D/J=0.15$, Opzione DM B), dove entrambe
le famiglie sono state ottimizzate con multistart a venti punti di
partenza casuali.

\begin{table}[htbp]
\centering
\small
\renewcommand{\arraystretch}{1.3}
\begin{tabular}{@{}lccc@{}}
\toprule
\rowcolor{grigetto}
\textbf{Ansatz} & \textbf{fidelity} $\mathcal F$ & \textbf{CNOT} & \textbf{gate a 1 qubit} \\
\midrule
RBS-2q.6 (6 par.) & $0.9999110739$ & $6$ & $46$ \\
$W$-2q.6 (6 par.) & $0.9999999999999996$ & $6$ & $25$ \\
\bottomrule
\end{tabular}
\caption{Fidelity sotto DM (Scenario A) delle due famiglie di ansatz a
sei parametri, dopo transpilazione in base $\{rz,\mathrm{sx},x,cx\}$ a
livello di ottimizzazione 3. Il numero di CNOT è identico (due per
blocco, tre blocchi); il numero di gate a un qubit di RBS è quasi il
doppio di quello di $W$, per via delle porte di Hadamard e delle due
rotazioni $R_y$ per blocco (contro l'unica di $W$).}
\label{tab:rbs-vs-w}
\end{table}
\FloatBarrier

Questo risultato è l'opposto di quanto accade nel dimero, dove RBS ---
più economico in gate --- è la scelta raccomandata e $W$ è più costoso
senza vantaggi. Nel trimero, RBS presenta un tetto strutturale reale,
seppur piccolo, mentre $W$ raggiunge l'esatto a precisione di macchina,
con un costo in gate inferiore: nessun compromesso da giustificare fra
le due metriche.

L'origine del tetto di RBS è geometrica, non un limite
dell'ottimizzatore. Il blocco RBS puro (senza le tre rotazioni $R_y$
finali che rompono la conservazione di $M$) resta confinato nel settore
di magnetizzazione a cui appartiene lo stato di partenza $X\ket{000}$,
cioè $M=+1/2$, Hamming weight $1$. Il peso della vera GS in quel settore,
a Scenario A, è $0.00098$; la fidelity massima raggiungibile da un
ansatz confinato in quel settore --- indipendentemente dal numero di
ripetizioni del blocco --- coincide con quel peso a sette cifre
($0.0009814687$, invariato ripetendo il blocco fino a tre giri). A
differenza del dimero, dove aggiungere due rotazioni locali indipendenti
dopo un solo blocco RBS ripara completamente il problema
($\mathcal F=1$), nel trimero le tre rotazioni $R_y$ finali riducono
drasticamente il tetto (da $\mathcal F\approx0.001$ a $\mathcal
F\approx0.9999$) ma non lo chiudono del tutto; aggiungere un secondo giro
completo di blocchi (nove parametri anziché sei) non sposta il risultato.
La differenza fra le due famiglie non è quindi solo una questione di
efficienza in gate, come nel dimero, ma anche di potere espressivo a
parità di parametri.

\subsection{Ansatz canonico}

Per il resto di questo lavoro si adotta l'ansatz \textbf{$W$-2q.6}: sei
parametri, fidelity esatta a precisione di macchina al punto Scenario A.

\begin{table}[htbp]
\centering
\small
\renewcommand{\arraystretch}{1.3}
\begin{tabular}{@{}lc@{}}
\toprule
\rowcolor{grigetto}
\textbf{Quantità} & \textbf{valore} \\
\midrule
$E_\text{VQE}$ & $-5.53437001739340$ \\
$E_\text{esatto}$ & $-5.53437001739350$ \\
$\Delta E$ & $9.8\times10^{-14}$ \\
$\mathcal F$ & $0.99999999999961$ \\
\bottomrule
\end{tabular}
\caption{Risultato dell'ansatz canonico $W$-2q.6 a Scenario A; multistart
a 12 punti di partenza casuali, tutti convergenti allo stesso ottimo.}
\end{table}
\FloatBarrier

Questo risultato è stato verificato non solo al singolo punto Scenario A,
ma su uno sweep di tredici valori di $b/J$ fra $0.3$ e $4.5$ (DM Opzione B
sempre acceso, $D/J=0.15$), inclusi punti ben sotto e ben sopra l'incrocio
e il punto dell'incrocio stesso ($b/J=2.372$): la fidelity resta
$\mathcal F=1.00000000$ su tutti i punti testati, con lo stesso stato
iniziale $X\ket{000}$ ovunque --- nessuna logica a rami di preparazione è
necessaria per questo ansatz, a differenza di quanto accade con ansatz
più minimali sul dimero.

\begin{figure}[htbp]
\centering
\begin{minipage}{0.48\textwidth}
\centering
\includegraphics[width=\textwidth]{figure/circ_rbs2q_anello}
\\ RBS-2q.6
\end{minipage}
\hfill
\begin{minipage}{0.48\textwidth}
\centering
\includegraphics[width=\textwidth]{figure/circ_w2q_anello}
\\ $W$-2q.6 (ansatz canonico)
\end{minipage}
\caption{Le due famiglie di ansatz in forma compatta: rotazione $X$ sul
primo qubit, un blocco a due qubit per ciascuno dei tre legami, tre
rotazioni $R_y$ finali indipendenti. Cambia solo il blocco a due qubit
usato sui legami.}
\label{fig:ansatz-completo}
\end{figure}
\FloatBarrier

\section{Evoluzione temporale: decomposizione di Suzuki--Trotter}
\label{sec:trotter}

\subsection{Una struttura a tre livelli}

Si scrive $H=H_0+H_\text{DM}$, con $H_0=bS_z^\text{tot}+H_\text{ex}$
(campo più scambio isotropo) e $H_\text{DM}$ il solo termine di
Dzyaloshinskii--Moriya. Un fatto esatto, non approssimato, semplifica il
quadro esterno: poiché $[S_z^\text{tot},H_\text{ex}]=0$,
\begin{equation}
e^{-iH_0\tau} = e^{-ib\tau S_z^\text{tot}}\, e^{-i\tau H_\text{ex}},
\qquad \text{esatto.}
\end{equation}
Fin qui, meglio del dimero: nessuna approssimazione a questo livello. Ma
con tre legami che condividono qubit --- il legame $(1,2)$ e il legame
$(2,3)$ condividono il qubit del sito $2$ --- lo scambio isotropo stesso
non si esponenzia in un solo passo:
\begin{equation}
[\vec\sigma_1\!\cdot\!\vec\sigma_2,\ \vec\sigma_2\!\cdot\!\vec\sigma_3] \neq 0.
\end{equation}
Serve dunque un Trotter \emph{interno} anche per $e^{-i\tau H_\text{ex}}$,
spezzato sui tre legami; lo stesso vale per $H_\text{DM}$ (Opzione B), che
condivide la stessa struttura a tre legami. Il risultato è una struttura a
tre livelli, uno in più rispetto al dimero:
\begin{itemize}[leftmargin=1.8em,itemsep=2pt]
\item \emph{livello 0} (esatto): $H_0=bS_z^\text{tot}+H_\text{ex}$
fattorizza in campo $\times$ scambio senza errore;
\item \emph{livello 1}: $e^{-i\tau H_\text{ex}}$, approssimato con un
singolo passaggio di Trotter sui tre legami;
\item \emph{livello 2}: $e^{-i\tau H_\text{DM}}$, stessa approssimazione,
stesso motivo.
\end{itemize}
Un Trotter esterno di primo ordine alterna poi $H_0$ e $H_\text{DM}$ su
$N$ passi, con errore $O((t/N)^2)$ per passo dovuto al solo fatto che
$[S_z^\text{tot},H_\text{DM}]\neq0$ (lo scambio non contribuisce
all'errore esterno, essendo esatto al livello 0).

\begin{figure}[htbp]
\centering
\includegraphics[width=\textwidth]{figure/circ_trotter_anello}
\caption{Un singolo passo esterno di Trotter. Nell'ordine: $e^{-i\tau
H_\text{ex}}$ (blocchi $R_{XX},R_{YY},R_{ZZ}$ per ciascuno dei tre
legami, livello 1), $e^{-i\tau bS_z^\text{tot}}$ (tre $R_z$, esatto),
$e^{-i\tau H_\text{DM}}$ (blocchi Hadamard--$R_{ZZ}$--Hadamard per
ciascun legame, livello 2). Il circuito completo è questo blocco ripetuto
$N$ volte.}
\label{fig:circ-trotter}
\end{figure}
\FloatBarrier

L'implementazione Qiskit del circuito a un passo è stata verificata
confrontando l'azione ripetuta $N$ volte del circuito con la potenza
$N$-esima dell'operatore corrispondente (\texttt{Operator(qc)}$^N$):
scarto massimo sulle ampiezze complesse $4.25\times10^{-15}$ a $t=6$,
$N=60$, coerente con la precisione di macchina attesa per un confronto
fra due modi di calcolare lo stesso operatore (non una verifica fisica
indipendente).

\subsection{Errore di Trotter e convergenza}

La dinamica è stata studiata al punto di lavoro $R_0$ ($b/J=0.05$,
$D/J=1.93$), scelto per dare un'evoluzione temporale non monocromatica
(più frequenze visibili simultaneamente), a partire dallo stato iniziale
$\ket{000}$.

\begin{figure}[htbp]
\centering
\includegraphics[width=0.8\textwidth]{figure/fig05_trotter_sz_anello}
\caption{$\langle S_z^\text{tot}\rangle(t)$ al punto $R_0$: evoluzione
esatta (nera) contro Trotter a $N=50,150,500$ passi.}
\label{fig:sz}
\end{figure}
\FloatBarrier

\begin{table}[htbp]
\centering
\small
\renewcommand{\arraystretch}{1.25}
\begin{tabular}{@{}rccc@{}}
\toprule
\rowcolor{grigetto}
$N$ & $\lVert U_\text{Trotter}-U_\text{esatto}\rVert$ & $1-\mathcal F$ & soglia raggiunta \\
\midrule
100  & $3.49\times10^{-1}$ & $5.03\times10^{-2}$ & --- \\
300  & $1.08\times10^{-1}$ & $4.90\times10^{-3}$ & $1-\mathcal F<10^{-2}$ \\
1000 & $3.15\times10^{-2}$ & $4.25\times10^{-4}$ & $1-\mathcal F<10^{-3}$ \\
3000 & $1.04\times10^{-2}$ & $4.68\times10^{-5}$ & $1-\mathcal F<10^{-4}$ \\
\bottomrule
\end{tabular}
\caption{Errore di Trotter a $t=6$, punto $R_0$. Le pendenze misurate in
scala doppio-logaritmica sono $-0.97$ (norma dell'operatore) e $-1.85$
($1-\mathcal F$), vicine ai valori teorici attesi $-1,-2$ ma non ancora
pienamente asintotiche a questi $N$.}
\label{tab:errore}
\end{table}
\FloatBarrier

Il regime pienamente asintotico richiede $N$ sensibilmente più grande che
nel dimero, dove $N=80$ era sufficiente per $1-\mathcal F\sim10^{-5}$: qui
servono centinaia--migliaia di passi per soglie comparabili. Questo non è
un peggioramento generico dovuto alla dimensione del sistema, ma è
imputabile specificamente al Trotter interno sui legami condivisi
(livelli 1 e 2), assente nel dimero, e a un rapporto $D/J\approx1.9$ non
piccolo al punto $R_0$.

\begin{figure}[htbp]
\centering
\includegraphics[width=0.68\textwidth]{figure/fig06_trotter_errore_anello}
\caption{Errore a $t=6$ (punto $R_0$), scala doppio-logaritmica. Le rette
punteggiate sono le pendenze asintotiche attese ($-1$ per la norma
dell'operatore, $-2$ per $1-\mathcal F$).}
\label{fig:errore}
\end{figure}
\FloatBarrier

\subsection{Costo in gate}

\begin{table}[htbp]
\centering
\small
\renewcommand{\arraystretch}{1.3}
\begin{tabular}{@{}lcc@{}}
\toprule
\rowcolor{grigetto}
\textbf{Compilazione} & \textbf{CNOT per passo} & \textbf{profondità} \\
\midrule
traduzione diretta, nessuna ottimizzazione & $30$ & $109$ \\
compilazione ottimizzata & $\mathbf{15}$ & $52$ \\
\bottomrule
\end{tabular}
\caption{Costo di un passo esterno di Trotter (tre legami, scambio più
DM), dopo transpilazione in base $\{rz,\mathrm{sx},x,cx\}$. La
compilazione ottimizzata dimezza il conteggio di CNOT rispetto alla
traduzione diretta.}
\end{table}
\FloatBarrier

Questo conteggio (15 CNOT per passo) è il dato che, nella seconda parte
di questo lavoro, determina il compromesso fra errore di Trotter ed
errore hardware sotto rumore.

\section{Correlatori dinamici}
\label{sec:correlatori}

\subsection{Definizione e circuito}

Si considera la stessa grandezza già definita nel caso del dimero,
estesa a tre siti:
\begin{equation}
C_{ij}^{\alpha\beta}(t) = \bra{\psi_0}\,\sigma_i^\alpha(t)\,\sigma_j^\beta(0)\,\ket{\psi_0},
\qquad \alpha,\beta\in\{x,y,z\},\quad i,j\in\{1,2,3\},
\end{equation}
per un totale di $81$ combinazioni. Poiché il prodotto
$\sigma_i^\alpha(t)\sigma_j^\beta(0)$ non è in generale hermitiano, la sua
misura richiede uno schema di tipo Hadamard test con un qubit ausiliario:
il circuito impiega quindi quattro qubit --- i tre del registro fisico
più l'ancilla, qui posta sul qubit $3$.

\begin{figure}[htbp]
\centering
\includegraphics[width=\textwidth]{figure/circ_correlatore_anello_re}
\caption{Circuito per la parte reale di $C_{21}^{xx}(t{=}1)$ (mostrato con
$N=1$ passo di Trotter per leggibilità; nell'uso effettivo, $N=60$).
Nell'ordine: preparazione dello stato fondamentale sul registro a tre
qubit; Hadamard sull'ancilla; $\sigma_1^x$ controllata; evoluzione $U(t)$
non controllata sul solo registro (il blocco di Trotter descritto in
\S\ref{sec:trotter}); $\sigma_2^x$ anti-controllata; rotazione di base
finale e misura dell'ancilla. La variante per la parte immaginaria
differisce solo nella rotazione finale ($R_x(\pi/2)$ anziché $H$).}
\label{fig:circ-correlatore}
\end{figure}
\FloatBarrier

Il circuito è stato verificato contro un riferimento classico calcolato
con implementazione indipendente (esponenziale di matrice diretto della
piena Hamiltoniana, nessun codice condiviso con il circuito): la
coincidenza è confermata su tredici istanti temporali, coerentemente con
l'errore di Trotter atteso a $N=60$.

\subsection{Zeri strutturali}

Al punto di lavoro Scenario A, delle $81$ combinazioni, $32$ risultano
nulle all'istante iniziale $t=0$: l'Hamiltoniana è reale e simmetrica,
quindi il fondamentale può scegliersi reale, e un prodotto
$\sigma_i^\alpha\sigma_j^\beta$ puramente immaginario ha media nulla su
uno stato reale. Di queste, solo \textbf{quattro} restano esattamente
nulle per ogni $t$: $C_{33}^{xz}$, $C_{33}^{yz}$, $C_{33}^{zx}$,
$C_{33}^{zy}$, tutte relative allo stesso sito (il sito $3$). Questo
sottoinsieme è stato verificato numericamente su una griglia di $81$
punti in $t\in[0,8]$; una spiegazione in forma chiusa di questi zeri
strutturali, specifici del sito 3 in questa parametrizzazione isoscele,
non è ancora stata derivata.

\begin{figure}[htbp]
\centering
\includegraphics[width=\textwidth]{figure/fig08_scan81_anello}
\caption{Tutte le 81 combinazioni al punto Scenario A. \textbf{Sinistra}:
valore all'istante iniziale. \textbf{Destra}: ampiezza massima su
$t\in[0,8]$. Delle $32$ caselle nulle a $t=0$, solo quattro restano nulle
su tutto l'intervallo.}
\label{fig:scan}
\end{figure}
\FloatBarrier

\subsection{Ricchezza spettrale}

Inserendo una risoluzione dell'identità sugli autostati di $H$, ciascun
correlatore si scompone in una somma di esponenziali,
$C_{ij}^{\alpha\beta}(t)=\sum_{n=0}^{7} c_n\,e^{-i(E_n-E_0)t}$. Il
criterio di ``ricchezza'' spettrale è quanto il peso $c_n$ sia distribuito
su più righe invece che concentrato su una sola. Due autocorrelazioni
sullo stesso sito possono avere contenuto spettrale molto diverso:
$C_{11}^{zz}$ distribuisce il peso su quattro righe comparabili (riga
dominante al $32\%$), mentre $C_{33}^{yy}$ lo concentra su due righe
vicine (riga dominante al $52\%$). Ampiezza e ricchezza spettrale sono
dunque proprietà indipendenti, non correlate al sito o alla coppia di siti
in sé, ma alle componenti di spin specifiche coinvolte.

\begin{figure}[htbp]
\centering
\includegraphics[width=\textwidth]{figure/fig09_ricco_piatto_anello}
\caption{Due correlatori a confronto, stessi assi. \textbf{Sinistra}
($C_{11}^{zz}$): segnale con battimento irregolare, più frequenze
visibili. \textbf{Destra} ($C_{33}^{yy}$): oscillazione più regolare,
quasi periodica.}
\end{figure}
\FloatBarrier

\begin{figure}[htbp]
\centering
\includegraphics[width=0.7\textwidth]{figure/fig10_spettro_righe_anello}
\caption{Peso delle transizioni disponibili (esclusa $n=0$, non
oscillante) per i due correlatori della figura precedente.}
\label{fig:spettro-righe}
\end{figure}
\FloatBarrier

\subsection{Effetto della preparazione VQE reale}

Tutti i risultati precedenti usano lo stato fondamentale esatto (ampiezze
dirette) come punto di partenza. Sostituendo questa preparazione con il
circuito $W$-2q.6 realmente eseguito (fidelity $\mathcal
F=0.9999999999996$ a Scenario A, praticamente esatta), l'effetto atteso è
piccolo. Sulla curva $C_{21}^{xx}(t)$ (13 istanti, $t\in[0,6]$), lo scarto
medio fra le due preparazioni è $2.50\times10^{-7}$, con massimo
$6.09\times10^{-7}$; ripetendo il confronto sull'intero scan delle 81
combinazioni a $t=2$, lo scarto medio è $2.90\times10^{-7}$, con massimo
$1.02\times10^{-6}$ (combinazione $C_{11}^{zz}$) e minimo
$6.04\times10^{-10}$.

\begin{figure}[htbp]
\centering
\includegraphics[width=\textwidth]{figure/fig_vqe_vs_esatto_correlatore}
\caption{$C_{21}^{xx}(t)$, preparazione esatta contro preparazione via
circuito $W$-2q.6, entrambe con $N=60$ passi di Trotter.
\textbf{Sinistra}: le due curve sono indistinguibili a questa scala.
\textbf{Destra}: lo stesso scostamento in scala logaritmica, ordine
$10^{-7}$.}
\label{fig:vqe-vs-esatto}
\end{figure}
\FloatBarrier

Poiché l'errore di preparazione è qui quasi sedici ordini di grandezza
più piccolo che nel caso del dimero con il suo ansatz meno accurato
($1-\mathcal F\sim4\times10^{-13}$ contro $1-\mathcal F\sim0.06$), lo
scarto propagato resta uniformemente minuscolo su tutte le combinazioni,
senza alcuna correlazione misurabile con l'ampiezza del correlatore
(coefficiente di correlazione $-0.02$ sulle 81 combinazioni): a questa
scala l'imperfezione di preparazione si comporta come rumore numerico
distribuito casualmente, non come un effetto strutturato.

Tutti i risultati di questa Parte sono ottenuti a simulazione ideale:
nessun rumore di gate, nessun errore di lettura. La seconda parte di
questo lavoro introduce questi effetti.

% =====================================================================
\part*{Parte II --- Il sistema aperto}
\addcontentsline{toc}{part}{Parte II --- Il sistema aperto}
% =====================================================================

\section{Il modello di rumore}
\label{sec:rumore}

Su hardware reale, ogni gate è affetto da un errore e la misura stessa
non è affidabile al cento per cento: il vettore di stato puro $\ket\psi$
va sostituito con l'operatore densità $\rho=\ket\psi\bra\psi$, e la sua
evoluzione sotto un canale di rumore $\mathcal N$ con
$\rho\to\mathcal N(\rho)$, in generale non più uno stato puro. Le
osservabili non cambiano forma, $\langle O\rangle=\Tr[\rho O]$, ma $\rho$
porta meno informazione di un vettore di stato: questo è il prezzo del
rumore.

Il modello adottato include due soli canali di rumore, non uno stack
generico di rumore hardware: un canale \textbf{depolarizzante di gate}, a
uno e due qubit, e un canale di \textbf{readout}. I processi $T_1,T_2$
sono esclusi dal modello. La calibrazione è basata su valori mediani
reali del dispositivo \texttt{ibm\_torino}:

\begin{table}[htbp]
\centering
\small
\renewcommand{\arraystretch}{1.3}
\begin{tabular}{@{}lc@{}}
\toprule
\rowcolor{grigetto}
\textbf{Parametro} & \textbf{valore (mediane \texttt{ibm\_torino})} \\
\midrule
$\varepsilon_{1q}$ (gate a 1 qubit) & $2.9\times10^{-4}$ \\
$\varepsilon_{2q}$ (gate a 2 qubit) & $3.8\times10^{-3}$ \\
$p_\text{readout}$ & $2.3\times10^{-2}$ \\
\bottomrule
\end{tabular}
\caption{Calibrazione del modello di rumore.}
\end{table}
\FloatBarrier

I parametri di contrazione della sfera di Bloch corrispondenti, per il
canale depolarizzante, sono $\lambda_{1q}=2\varepsilon_{1q}=5.80\times10^{-4}$
e $\lambda_{2q}=\tfrac43\varepsilon_{2q}=5.07\times10^{-3}$.

\begin{figure}[htbp]
\centering
\includegraphics[width=0.55\textwidth]{figure/fig_depolarizzazione_bloch_anello}
\caption{Contrazione della sfera di Bloch sotto il canale depolarizzante,
$1-\lambda$ in funzione di $\lambda$. Le linee verticali segnano i
valori di calibrazione $\lambda_{1q}$ e $\lambda_{2q}$.}
\end{figure}
\FloatBarrier

Il modello di rumore non dipende dal numero di qubit del circuito a cui
viene applicato: la costruzione dell'oggetto \texttt{NoiseModel} in
Qiskit tramite \texttt{add\_all\_qubit\_quantum\_error} è per costruzione
indipendente dalla topologia del registro. Il rumore di lettura è
trattato come simmetrico in questa e nelle successive sezioni, salvo
laddove esplicitamente generalizzato (\S\ref{sec:readout-asimmetrico}).
Non sono inclusi modelli di crosstalk fra qubit vicini, né dipendenza
spaziale del rumore: si assume lo stesso $\varepsilon$ su ogni qubit e
coppia di qubit.

\section{Il VQE sotto rumore}
\label{sec:vqe-rumore}

I parametri ottimali $\theta^*$ dell'ansatz $W$-2q.6, determinati senza
rumore al punto Scenario A ($b/J=2.4$, $D/J=0.15$, Opzione DM B), vengono
qui eseguiti su un simulatore rumoroso, senza alcuna riottimizzazione: si
misura così la sola degradazione di un risultato già ideale. Le
grandezze rilevanti sono $E_\text{rumoroso}=\Tr[\rho H]$ e
$F_\text{rumorosa}=\bra{\psi_0}\rho\ket{\psi_0}$.

\begin{table}[htbp]
\centering
\small
\renewcommand{\arraystretch}{1.3}
\begin{tabular}{@{}lc@{}}
\toprule
\rowcolor{grigetto}
\textbf{Quantità} & \textbf{valore} \\
\midrule
$E$ (rumoroso, calibrazione \texttt{ibm\_torino}) & $-5.37687489$ \\
$E_\text{esatto}$ & $-5.53437002$ \\
$F$ (rumorosa) & $0.97021471$ \\
$1-F$ & $2.98\%$ \\
CNOT nel circuito & $6$ \\
\bottomrule
\end{tabular}
\caption{VQE sotto rumore, Scenario A, parametri ideali riusati senza
riottimizzazione.}
\end{table}
\FloatBarrier

La perdita di fedeltà, $1-F=2.98\%$, è circa il doppio di quella
corrispondente nel dimero a parità di calibrazione di rumore relativa.
La ragione è strutturale: l'ansatz del dimero impiega $2$ CNOT, contro i
$6$ di questo ansatz --- tre volte tanti canali di rumore attraversati,
per un ansatz che è anche il più costoso in gate dell'intera pipeline del
trimero, non il più economico come nel caso a due spin. La dipendenza
dell'energia e della fedeltà rumorose da $\varepsilon_{2q}$ è monotona,
come nel dimero, ma su una scala verticale più ampia a parità di
$\varepsilon_{2q}$.

\section{La dinamica sotto rumore}
\label{sec:dinamica-rumore}

Nel sistema chiuso (\S\ref{sec:trotter}), l'errore di Trotter scende
monotonamente con $N$: più passi, meglio è, l'unico costo essendo la
profondità del circuito. Sotto rumore questo non è più vero, perché ogni
passo aggiunge $15$ CNOT rumorosi (il costo determinato in
\S\ref{sec:trotter}): esiste un valore ottimale $N^*$ che massimizza la
fedeltà rumorosa, né il più piccolo né il più grande possibile.

Sono stati trattati due scenari, che isolano contributi diversi
all'errore complessivo:
\begin{itemize}[leftmargin=1.8em,itemsep=3pt]
\item \textbf{Scenario A} ($b/J=2.4$, $D/J=0.15$): preparazione VQE
($W$-2q.6) seguita da $N$ passi di Trotter --- la curva $F(N)$ riflette
qui due sorgenti di errore sovrapposte, l'imperfezione della preparazione
rumorosa (\S\ref{sec:vqe-rumore}, $F\approx0.970$) e l'errore di Trotter
più rumore di gate sul circuito che segue.
\item \textbf{Scenario B} (punto $R_0$: $b/J=0.05$, $D/J=1.93$): nessuna
preparazione, si parte direttamente da $\ket{000}$ --- zero rumore di
preparazione. La curva $F(N)$ mostra qui solo l'effetto di Trotter più
rumore di gate, isolato dal contributo della preparazione.
\end{itemize}

\begin{figure}[htbp]
\centering
\includegraphics[width=\textwidth]{figure/fig_passo3_F_vs_N_anello}
\caption{Fedeltà rumorosa $F(N)$ rispetto allo stato bersaglio fisico, i
due scenari a confronto. \textbf{Sinistra}: Scenario A, $N^*=3$.
\textbf{Destra}: Scenario B, $N^*=9$, con un minimo locale marcato
attorno a $N=2$--$3$ prima di risalire al massimo vero.}
\label{fig:FvsN}
\end{figure}
\FloatBarrier

I valori ottenuti sono $N^*_A=3$, $F(N^*)=0.700335$, e $N^*_B=9$,
$F(N^*)=0.381896$. Il minimo locale della curva B, attorno a $N=2$, è la
stessa fisica del ``minimo non perturbativo'' già osservata nel dimero:
a pochi passi, l'errore di Trotter può temporaneamente aumentare prima
che il guadagno di accuratezza prenda il sopravvento --- una proprietà
della dinamica stessa, non un artefatto del rumore. Nessuna misura è
coinvolta in questa fedeltà, calcolata direttamente sull'operatore
densità.

\section{I correlatori sotto rumore}
\label{sec:correlatori-rumore}

A differenza del VQE e della dinamica, il correlatore dinamico richiede
una misura reale dell'ancilla, e quindi introduce anche il rumore di
lettura, non solo quello di gate. Il readout è implementato come misura
genuinamente campionata: un \texttt{ReadoutError} vero agganciato al
simulatore, con $10^5$ shot per punto, senza alcuna formula di
correzione analitica applicata a posteriori.

Il correlatore di riferimento è $C_{21}^{xx}(t{=}2)$, al punto Scenario
A --- una combinazione rappresentativa (non l'estremo più debole del
segnale fra le 81 disponibili).

\begin{figure}[htbp]
\centering
\includegraphics[width=0.65\textwidth]{figure/fig_passo4_correlatore_vs_N_anello}
\caption{$|C_{21}^{xx}(t{=}2,N)|$, rumore di riferimento, readout a shot
finiti ($10^5$ shot per punto). $N^*=3$, $|C(N^*)|=0.5026$.}
\end{figure}
\FloatBarrier

Un confronto a tre livelli --- valore esatto, Trotter senza rumore,
Trotter con rumore completo --- permette di quantificare separatamente
l'effetto della discretizzazione e quello del rumore:

\begin{table}[htbp]
\centering
\small
\renewcommand{\arraystretch}{1.3}
\begin{tabular}{@{}lccc@{}}
\toprule
\rowcolor{grigetto}
\textbf{Livello} & $N^*$ & $|C|$ & \textbf{rapporto a $C_\text{esatto}$} \\
\midrule
$C_\text{esatto}(t{=}2)$ --- nessun Trotter, nessun rumore & --- & $0.6723$ & $1.00\times$ \\
$C(N;t{=}2)$, rumore nullo (solo Trotter) & $3$ & $0.6773$ & $1.007\times$ \\
$C(N;t{=}2)$, rumore vero (Trotter $+$ gate $+$ lettura) & $3$ & $0.5026$ & $0.748\times$ \\
\bottomrule
\end{tabular}
\caption{I tre livelli per il correlatore $C_{21}^{xx}$, Scenario A.
``Esatto'' indica il calcolo numerico senza approssimazioni fisiche
(diagonalizzazione diretta più esponenziale di matrice diretto), non una
formula chiusa.}
\end{table}
\FloatBarrier

Un punto rilevante per la generalità del criterio $N^*=\arg\max_N|C(N)|$:
a rumore nullo, il valore $|C(N^*)|$ è qui praticamente coincidente con
il valore esatto ($1.007\times$). Questo risultato non va assunto per
analogia con altri sistemi: il criterio $N^*=\arg\max_N|C(N)|$ massimizza
il segnale misurabile, non necessariamente l'accuratezza della stima
rispetto al valore fisico vero, e il grado di sovra- o sotto-stima
introdotto dipende dalla combinazione specifica di errore di Trotter e
forma della curva $|C(N)|$ per il sistema e il punto di lavoro
considerati; non è quindi una proprietà universale del criterio, e va
verificato caso per caso.

\section{Scan sui parametri di rumore}
\label{sec:scan-rumore}

Per valutare quanto i risultati precedenti dipendano dalla calibrazione
specifica di \texttt{ibm\_torino}, $N^*$ è stato studiato al variare dei
parametri di rumore, uno alla volta, con gli altri fissati al valore di
riferimento.

\begin{figure}[htbp]
\centering
\includegraphics[width=\textwidth]{figure/fig_passo5_Nstar_vs_eps_anello}
\caption{$N^*$ in funzione di $\varepsilon_{2q}$ (scala logaritmica).
\textbf{Sinistra}: fedeltà di Trotter, Scenario A. \textbf{Destra}:
correlatore, readout a shot finiti.}
\end{figure}
\FloatBarrier

\begin{table}[htbp]
\centering
\small
\renewcommand{\arraystretch}{1.3}
\begin{tabular}{@{}lcc@{}}
\toprule
\rowcolor{grigetto}
$\varepsilon_{2q}$ & $N^*$ Trotter & $N^*$ correlatore \\
\midrule
$1\times10^{-3}$ & 4 & 3 \\
$2\times10^{-3}$ & 4 & 3 \\
$3.8\times10^{-3}$ (riferimento) & 3 & 3 \\
$6\times10^{-3}$ & 2 & 3 \\
$1\times10^{-2}$ & 2 & 3 \\
\bottomrule
\end{tabular}
\caption{$N^*$ Trotter è monotono non-crescente in $\varepsilon_{2q}$,
come atteso (più rumore, meno conviene allungare il circuito); $N^*$
correlatore resta invece fermo a $3$ su tutto il range esplorato, una
differenza fra le due metriche non ancora indagata in dettaglio.}
\end{table}
\FloatBarrier

L'invarianza di $N^*$ del correlatore rispetto a $p_\text{readout}$ è
stata verificata esplicitamente a shot finiti su tre valori
($p=0.00,\,0.05,\,0.20$), sempre con $N^*=3$: risultato atteso
analiticamente (la correzione di readout simmetrico è un fattore
moltiplicativo costante in $N$, che non può spostare la posizione di un
massimo) e qui confermato anche con una misura genuinamente campionata,
non solo assunto dall'argomento analitico.

Con la conclusione di questo scan si chiudono i cinque stadi principali
della pipeline sotto rumore. Le due sezioni seguenti trattano due
estensioni, motivate da domande lasciate aperte: se convenga
riottimizzare il VQE vedendo direttamente il rumore, e cosa cambi quando
il readout smette di essere simmetrico.

\section{Estensione: VQE noise-aware}
\label{sec:noise-aware}

Per un canale depolarizzante isotropo, l'energia rumorosa si scrive come
trasformazione affine dell'energia ideale:
\begin{equation}
E_\text{rumoroso}(\boldsymbol\theta) = (1-\lambda)\,E_\text{ideale}(\boldsymbol\theta)
+ \lambda\,\frac{\Tr[H]}{d},
\end{equation}
dove $d$ è la dimensione dello spazio di Hilbert. Poiché la
trasformazione è affine con coefficiente positivo, lo stesso
$\boldsymbol\theta^*$ che minimizza $E_\text{ideale}$ minimizza anche
$E_\text{rumoroso}$: un argomento generale, indipendente dal sistema, che
predice nessun vantaggio a riottimizzare vedendo il rumore durante
l'ottimizzazione stessa (VQE noise-aware) rispetto al riuso dei parametri
ideali.

Questo argomento presuppone però una struttura di circuito fissa --- lo
stesso numero di gate rumorosi per qualunque $\boldsymbol\theta$. Con un
ansatz a sei parametri, il transpilatore può semplificare il blocco
CNOT--$R_y(\theta)$--CNOT a valori speciali di $\theta$ (vicini a $0$ o
$\pi$), riducendo il conteggio di CNOT rispetto al caso generico; questo
non accade con un ansatz a struttura fissa (livello di ottimizzazione 0
della transpilazione), dove il conteggio resta costante qualunque sia
$\theta$, incluso ai valori speciali. Per isolare l'effetto fisico da
questo comportamento del compilatore, i risultati principali di questa
sezione sono ottenuti a struttura di circuito fissa.

\begin{table}[htbp]
\centering
\small
\renewcommand{\arraystretch}{1.3}
\begin{tabular}{@{}lcc@{}}
\toprule
\rowcolor{grigetto}
\textbf{Quantità} & \textbf{riuso $\theta^*_\text{ideale}$} & \textbf{VQE noise-aware} \\
\midrule
$E$ (rumoroso) & $-5.36232432$ & $-5.36648903$ \\
$F$ (rumorosa) & $0.96728547$ & $0.95153755$ \\
CNOT & 6 & 6 \\
\bottomrule
\end{tabular}
\caption{Confronto fra riuso dei parametri ideali e riottimizzazione
sotto rumore (VQE noise-aware), Scenario A, struttura di circuito
fissa, multistart a 12 punti casuali più il seme ideale.}
\end{table}
\FloatBarrier

Il risultato, $\Delta E=-4.16\times10^{-3}$ (energia migliore) e $\Delta
F=-1.57\times10^{-2}$ (fedeltà peggiore), mostra che la riottimizzazione
trova genuinamente un punto diverso da $\theta^*_\text{ideale}$: qui
l'argomento affine garantisce solo che $\theta^*_\text{ideale}$ resti un
minimo di $E_\text{ideale}$, non che sia l'unico minimo rilevante di
$E_\text{rumoroso}$ quando la dimensionalità dei parametri cresce.

A struttura fissa, tre semi di partenza convergono a un ottimo identico
entro $2\times10^{-11}$: il risultato è stabile.

A valle, la fedeltà di Trotter (Scenario A) e il correlatore dinamico
mostrano entrambi $N^*=3$, invariato rispetto al riuso dei parametri
ideali, verificato in entrambi i casi a shot finiti per il correlatore.

\begin{figure}[htbp]
\centering
\includegraphics[width=0.72\textwidth]{figure/fig_F_vs_N_confronto_trimero_anello}
\caption{Fedeltà di Trotter, struttura fissa. $N^*=3$ in entrambi i
casi; scostamento piccolo, su scala $\sim10^{-2}$.}
\end{figure}
\FloatBarrier

L'invarianza di $N^*$ qui non si spiega, come nel dimero, con un
argomento di continuità (lì $\theta^*_\text{noise-aware}$ era quasi
indistinguibile da $\theta^*_\text{ideale}$): in questo caso lo
spostamento dei parametri non è piccolo, essendo $\Delta E$ e $\Delta F$
tre ordini di grandezza più grandi che nel dimero. L'invarianza osservata
riflette piuttosto un'insensibilità dell'osservabile a valle a questo
particolare spostamento di parametri, non una proprietà generale che si
estenda automaticamente ad altri punti di lavoro.

\section{Estensione: readout asimmetrico}
\label{sec:readout-asimmetrico}

Il readout simmetrico ($p_{01}=p_{10}=p$) adottato nelle sezioni
precedenti è una semplificazione: su hardware reale il rilassamento
$T_1$ durante la misura rende tipicamente $p_{10}>p_{01}$. Con una
matrice di confusione generale,
\begin{equation}
\langle Z\rangle_\text{readout} = \alpha\langle Z\rangle_\text{ideale}+\beta,
\qquad \alpha=1-p_{01}-p_{10},\quad \beta=p_{10}-p_{01},
\end{equation}
per il correlatore complesso questo corrisponde a una traslazione
costante nel piano complesso, che non è garantita a priori preservare
l'invarianza di $N^*=\arg\max_N|C(N)|$ (somma e modulo non commutano).

Il circuito dei correlatori misura un solo bit classico, sempre e solo
l'ancilla: la correzione di readout asimmetrico si applica a quella
singola misura, indipendentemente dalla dimensione del registro fisico
--- il registro a tre qubit del trimero, più grande di quello a due qubit
del dimero, non introduce quindi alcuna complicazione aggiuntiva per
questa correzione.

Con uno split illustrativo $p_{10}/p_{01}=3\times$ (stessa media
$p_\text{readout}=2.3\times10^{-2}$: $p_{01}=0.0115$, $p_{10}=0.0345$),
tutto a shot finiti:

\begin{figure}[htbp]
\centering
\includegraphics[width=0.62\textwidth]{figure/fig_correlatore_N_asimmetrico_anello}
\caption{$|C_{21}^{xx}(t{=}2,N)|$, readout simmetrico contro asimmetrico
($8\times10^4$ shot per punto). $N^*=3$ in entrambi i casi.}
\end{figure}
\FloatBarrier

\begin{table}[htbp]
\centering
\small
\renewcommand{\arraystretch}{1.3}
\begin{tabular}{@{}lcc@{}}
\toprule
\rowcolor{grigetto}
& $N^*$ & $|C(N^*)|$ \\
\midrule
simmetrico ($p_{01}=p_{10}=0.023$) & 3 & $0.5059$ \\
asimmetrico ($3\times$) & 3 & $0.4934$ \\
\bottomrule
\end{tabular}
\end{table}
\FloatBarrier

Uno stress test oltre il range realistico, con rapporti
$p_{10}/p_{01}\in\{1,3,5,10,20,50\}$ (ben oltre il range tipicamente
osservato su hardware IBM, $2$--$3.5\times$), dà $N^*=3$ su tutti i
rapporti testati, verificato a shot finiti su ciascun punto. La
spiegazione è che $N^*$ è definito dal massimo di una curva $|C(N)|$ già
ripida attorno al suo picco: una traslazione costante sposta la curva
intera, ma raramente abbastanza da spostare la posizione del suo
massimo, finché $\alpha>0$ (cioè finché $p_{01}+p_{10}<1$, sempre vero
nel range fisico qui esplorato).

\section{Conclusioni}

Il passaggio dal dimero al trimero ad anello introduce, oltre alla
frustrazione geometrica prevista, una serie di conseguenze tecniche
concrete che modificano in modo non banale le conclusioni già stabilite
sul sistema più semplice:

\begin{itemize}[leftmargin=1.6em,itemsep=3pt]
\item il termine di Dzyaloshinskii--Moriya ammette più forme possibili,
e solo una di esse (Opzione B) apre l'incrocio protetto dalla
conservazione di $S_{12}^2$;
\item nel confronto fra ansatz variazionali, l'ordine di preferenza
osservato nel dimero si inverte: $W$-2q.6 domina RBS-2q sia in costo sia
in potere espressivo, mentre nel dimero era vero il contrario;
\item l'evoluzione temporale richiede un livello aggiuntivo di
approssimazione di Trotter, dovuto alla condivisione di qubit fra legami,
con una convergenza sensibilmente più lenta;
\item sotto rumore, la maggiore profondità dell'ansatz canonico si
traduce in una degradazione di fedeltà doppia rispetto al dimero a parità
di calibrazione; il criterio $N^*$ per il correlatore dinamico, a
differenza del dimero, non sovrastima qui il valore fisico vero; e la
riottimizzazione noise-aware del VQE produce, a differenza del dimero, un
compromesso energia/fedeltà misurabile, sebbene senza conseguenze
sull'$N^*$ osservato a valle.
\end{itemize}

Questi risultati indicano che le conclusioni metodologiche stabilite su
un sistema a due spin non si trasferiscono automaticamente a un sistema
di dimensione anche di poco superiore, e che ciascun confronto --- forma
del termine di interazione, scelta dell'ansatz, comportamento del
criterio $N^*$ --- richiede una verifica specifica per il sistema e il
punto di lavoro considerati.

\bigskip
\noindent\footnotesize\textit{Riferimenti.}\\
K.~Kambe, \emph{On the Paramagnetic Susceptibilities of Some Polynuclear
Complex Salts}, J.\ Phys.\ Soc.\ Jpn.\ \textbf{5}, 48 (1950).\\
A.~Crippa et al., \emph{Simulating Static and Dynamic Properties of
Magnetic Molecules with Prototype Quantum Computers}, Magnetochemistry
\textbf{7}, 117 (2021).\\
N.~Hatano, M.~Suzuki, \emph{Finding Exponential Product Formulas of
Higher Orders}, in \emph{Quantum Annealing and Other Optimization
Methods}, Lecture Notes in Physics \textbf{679}, Springer (2005).

\end{document}
